\section{System Architecture and Stage Interfaces}
\label{sec:architecture}

Our system is a two-stage cascade: the object detector of
Sec.~\ref{sec:detector} first localizes the animal and decides its species,
and a species-specific fine-grained classifier then assigns the breed. Images
in which no target animal is detected are rejected ($-1$) without ever
reaching a classifier. This detect-then-classify design follows a long line
of fine-grained recognition work, in which localizing the object before
classification removes background clutter and lets the classifier spend its
capacity on discriminative details \cite{zhang2014part}. The same cascade
pattern is deployed at scale in wildlife monitoring, where MegaDetector
\cite{megadetector2026} localizes animals and SpeciesNet \cite{speciesnet2025}
classifies the resulting crops; as discussed in Sec.~\ref{sec:detector}, we
evaluated that stack directly before settling on the lighter YOLOv6-based
\cite{li2022yolov6} realization of the same idea, and the runtime analysis of
Sec.~\ref{sec:speed} confirms the choice: even our lightweight detector
dominates end-to-end latency at ${\sim}83\%$ of the per-image budget, so a
heavier localization stage would have threatened the 5-second constraint.

The two stages communicate through a single, deliberately narrow interface:
the detector emits a crop and a species tag (\texttt{cat}, \texttt{dog}, or
\texttt{neither}), and the tag routes the crop to one of two independent
classifiers. Keeping this interface minimal was a conscious design decision
with three architectural consequences. First, each classifier solves a
10-class problem over visually related breeds instead of a 20-class problem
spanning two species; its entire output space is conditioned on the
detector's decision, so no capacity is spent on inter-species separation that
the detector already performs reliably. Second, rejection becomes a property
of the detector rather than an extra class the classifiers must learn:
confounder images produce no cat/dog detection above threshold and are
rejected before classification, so the classifiers never require training
data for a ``none of the above'' concept---an important simplification given
how our dataset was assembled from breed-labeled sources
(Sec.~\ref{sec:dataset}) that contain no negative examples. Third, the two
classifiers can be developed, retrained, and replaced independently, which in
practice allowed the training tracks described in Sec.~\ref{sec:training} to
iterate on the cat model---the harder of the two---without touching the dog
model at all.

The cost of species routing is that a routing mistake is unrecoverable: the
wrong-species model does not contain the correct breed in its output space.
In our end-to-end evaluation such cross-species failures were rare and
concentrated on atypical animals---the canonical example being a hairless
Sphynx cat detected as a dog, whereupon the dog model, forced to answer,
selected the wrinkle-faced Pug. We accepted this bounded cost in exchange for
the simpler, better-matched per-species problems; the performance asymmetry
between the two species models that this split exposed is analyzed in
Sec.~\ref{sec:training}.
